% Generated draft: requires mathematical and source review.
\documentclass[11pt]{article}
\usepackage[margin=25mm]{geometry}
\usepackage{fontspec}
\setmainfont{Latin Modern Roman}
\usepackage{amsmath,amssymb,mathtools,unicode-math}
\setmathfont{Latin Modern Math}
\usepackage{enumitem,needspace,xurl,hyperref}
\setlength{\parindent}{0pt}
\setlength{\parskip}{5pt}
\setlength{\emergencystretch}{3em}
\newcommand{\N}{\mathbb{N}}
\newcommand{\Z}{\mathbb{Z}}
\newcommand{\Q}{\mathbb{Q}}
\newcommand{\R}{\mathbb{R}}
\newcommand{\C}{\mathbb{C}}
\newcommand{\F}{\mathbb{F}}
\newcommand{\A}{\mathbb{A}}
\newcommand{\PP}{\mathbb{P}}
\newcommand{\E}{\mathbb{E}}
\newcommand{\eps}{\varepsilon}
\newcommand{\dd}{\,\mathrm{d}}
\begin{document}
\section*{UnsolvedMath ID: 30005231}
\textbf{Source problem code:} catalogue entry
\subsection{Definitions and mathematical statement}
The supplied record concerns the Ising model and the random-cluster model, the latter described as dependent edge percolation with parameters \((p,q)\). For integer \(q\), the excerpt states its connection with the \(q\)-state Potts model, including the Ising case \(q=2\), using the convention \(p=1-e^{-\beta}\). The underlying graphs, parameter ranges, boundary conditions, and dynamics needed to specify the intended results are not supplied completely.

One possible meaning of monotonicity for a Markov process is the following. On a partially ordered measurable state space \((\Omega,\preceq)\), a transition semigroup \((P_t)_{t\geq0}\) is order preserving if, for every bounded measurable increasing function \(f\), every \(t\geq0\), and every \(x,y\in\Omega\),
\[
 x\preceq y\quad\Longrightarrow\quad (P_tf)(x)\leq(P_tf)(y).
\]
Here increasing means that \(f(x)\leq f(y)\) whenever \(x\preceq y\). This definition is explanatory only: the supplied fragment does not establish whether the intended assumptions are this dynamical property, stochastic monotonicity of equilibrium measures in parameters or boundary conditions, or another monotonicity property.

\textbf{Recoverable question.} Can the monotonicity assumptions used in the Ising and random-cluster interacting-particle-system results referred to by the source be removed while retaining the intended conclusions under the other intended hypotheses?

\textbf{Unresolved scope.} This is a question fragment, not a fully specified quantified conjecture. The phrase referring to the preceding results has no complete antecedent in the supplied record. Neither the particular theorems nor their hypotheses, conclusions, replacement model class, and quantifiers can be recovered. Thus the question above cannot faithfully be strengthened into a precise theorem statement without additional source material.

\textbf{Supplied contextual claims.} The excerpt discusses random-cluster Glauber dynamics and its connection with the Swendsen--Wang algorithm. It reports rapid mixing from worst-case initialization on \(T_n\) for \(p<p_c(q,d)\), and describes rapid mixing for all \(p>p_c(q,d)\) as expected, with results reported for \(d=2\) and for sufficiently large \(p\), written \(p\gg p_c(q,d)\). It places phase coexistence and slow mixing at \(p=p_c(q,d)\) for large \(q\), describing a wired phase with a giant connected component and a free phase with microscopic components. These are contextual reports, not independently verified assertions or a specification of the target theorem. The excerpt does not define \(T_n\), specify the mixing-time bounds or boundary conditions, or quantify the large-parameter qualifications. Its final sentence, concerning analogues in reference [8] of results in reference [7], ends at ``microscopic windo'' and supplies no complete window or conclusion.
\subsection{Short English statement}
The record asks whether results for Ising and random-cluster dynamics can survive removal of their monotonicity assumptions. The relevant results and the exact meaning of monotonicity are missing, so a fully quantified mathematical conjecture cannot be recovered. The supplied open status and literature assessment are dated curation requiring independent verification.
\subsection{Sources}
\begin{itemize}
\item Supplied local UnsolvedMath registry record, permanent ID \(30005231\), source problem code \texttt{catalogue entry}, entitled ``Removing Monotonicity Assumptions in Interacting Particle Systems,'' in the Oberwolfach Reports Open Problems collection. No proposer is supplied; the proposed year is recorded as 2022.
\item Supplied report citation: \emph{Large Scale Stochastic Dynamics (2023)}, Oberwolfach Reports, DOI 10.4171/owr/2022/41, pp. 2399--2466. \url{https://doi.org/10.4171/owr/2022/41}. The supplied report title year is 2023, whereas the DOI and proposed-year metadata contain 2022; this discrepancy is unresolved here. The report is described as documenting workshop discussions and open problems. The supplied excerpt is truncated and does not identify a speaker or the intended preceding theorems.
\item Supplied literature source: Candellero--Stauffer, ``non-monotonicity.'' \url{https://kclpure.kcl.ac.uk/portal/en/publications/first-passage-percolation-in-hostile-environment-is-not-monotone-2}.
\item Supplied literature source: Gheissari, ``random-cluster dynamics.'' \url{https://doi.org/10.1002/rsa.21191}.
\item Literature assessment dated exactly \texttt{2026-08-22}: supplied dated curation, not independently verified. It labels the problem open and describes it as an unresolved fragment about essential monotonicity assumptions, asserting that subsequent work still treats relevant hostile-environment models as non-monotone rather than providing a general resolution. The relationship of those models to the intended Ising and random-cluster results is not established by the supplied fragment.
\item The excerpt cites numbered references [1], [3], [6], [7], [8], [9], and [11], but their bibliographic entries are absent. The registry labels its statement \texttt{corrected\_verified} and says only formula delimiters, notation, grammar, and typography were repaired with scope preserved. These supplied labels do not constitute independent mathematical or source verification.
\end{itemize}
\end{document}
